% Compile with: uplatex file.tex && dvipdfmx file.dvi
%           or: latexmk -uplatex file.tex
%\documentclass[uplatex,dvipdfmx]{jsarticle}
\documentclass[dvipdfmx,a4paper,11pt]{jsarticle}
\usepackage{amsmath, amssymb}
\usepackage{tikz}

\title{scenario\_jaynes\_cummings}
\date{\today}

\begin{document}
\maketitle

Jaynes-Cummings モデル（光子と2準位分子の相互作用）を扱う。光子の周波数 $\omega$、分子の遷移周波数 $\omega_0$、光・物質結合定数 $g$ を実数パラメータとして準備する。

omega を実数パラメータとして定義して

\[
  \omega
\]

$\omega_0$ を実数パラメータとして定義して

\[
  \omega_{0}
\]

g を実数パラメータとして定義して

\[
  g
\]

光子の消滅・生成演算子 $a,a^\dagger$ をボゾン演算子として、分子の2準位（スピン1/2）演算子 $S$ をスピン演算子として準備する。

a を光子の生成演算子とする

\[
  a
\]

S をスピン演算子として定義して

\[
  S, S^{x}, S^{y}, S^{z}
\]

Jaynes-Cummings ハミルトニアン $H=\omega a^\dagger a+\omega_0S_z+g(aS^++a^\dagger S^-)$ を、$S^\pm=S_x\pm iS_y$ として展開した形で登録する。第1項は光子のエネルギー、第2項は分子の励起エネルギー、第3項は光子の吸収・放出による光・物質相互作用を表す。

$\omega a'(1) a(1) + \omega_0 Sz(1) + g (a(1) (Sx(1) + i Sy(1)) + a'(1) (Sx(1) - i Sy(1)))$

\begin{equation}
  \omega a^{\dagger}(1) a(1) + \omega_0 Sz(1) + g (a(1) (Sx(1) + i Sy(1)) + a^{\dagger}(1) (Sx(1) - i Sy(1)))
  \label{eq:E1}
\end{equation}

光子演算子 $a$ のハイゼンベルク方程式（運動方程式）を求める。

a(1) の Heisenberg 方程式 E1

\begin{equation}
  \frac{d a(1)}{d t} = - g \mathit{i} Sx_1 - g Sy_1 - \omega \mathit{i} {a}_{1}
  \label{eq:E2}
\end{equation}

分子側の演算子 $S_x$ のハイゼンベルク方程式を求める。

Sx(1) の Heisenberg 方程式 E1

\begin{equation}
  \frac{d Sx(1)}{d t} = - 2 \omega_{0} Sy_1 - 2 g \mathit{i} {a}_{1}^{\dagger} Sz_1 + 2 g \mathit{i} Sz_1 {a}_{1}
  \label{eq:E3}
\end{equation}

励起数演算子 $N_{\rm ex}=a^\dagger a+S_z/2$（光子数と分子のスピン $z$ 成分の和）を定義する。

a'(1)*a(1) + Sz(1)/2

\begin{equation}
  a^{\dagger}(1) a(1) + Sz(1)/2
  \label{eq:E4}
\end{equation}

ハミルトニアン $H$ と励起数演算子 $N_{\rm ex}$ の交換子を計算する。物理的には、$H$ の3項はそれぞれ単独で全励起数を保存する——光子エネルギー項と分子エネルギー項は励起数基底で対角であり、相互作用項 $g(aS^++a^\dagger S^-)$ は光子1個を分子励起1個に変換する（またはその逆）だけで、$a^\dagger a+S_z$ の合計は変化しない（光子の吸収・放出）。CAS が計算した交換子がちょうど0に潰れるのはこのためであり、決め打ちの答えではなく、実際に計算された打ち消し合いである。（CAS は演算子の恒等式を適用して最終的な簡約値を直接返すため、項ごとの展開過程は別ステップとしては表示されない。）

[E1, E4]

\begin{equation}
  0
  \label{eq:E5}
\end{equation}

\end{document}